% Ad Atticum 6.4 (ad-atticum-06-04)
% Source: https://raw.githubusercontent.com/PerseusDL/canonical-latinLit/master/data/phi0474/phi057/phi0474.phi057.perseus-lat1.xml
% Fetched by scripts/fetch_latin.py

% Dateline (Perseus): Scr. in itinere paulo post Non. Iun. a. 704 (50) .

\ciceroLetterOpener{CICERO ATTICO salutem}

\ciceroSection{1} Tarsum venimus Nonis Iuniis. ibi me multa moverunt, magnum in Syria bellum, magna in Cilicia latrocinia, mihi difficilis ratio administrandi, quod paucos dies habebam reliquos annui muneris, illud autem difficillimum, relinquendus erat ex senatus consulto qui praeesset. nihil minus probari poterat quam quaestor Mescinius. nam de Coelio nihil audiebamus. rectissimum videbatur fratrem cum imperio relinquere; in quo multa molesta, discessus noster, belli periculum, militum improbitas, sescenta praeterea. o rem totam odiosam! sed haec fortuna viderit, quoniam consilio non multum uti licet.

\ciceroSection{2} tu quando Romam salvus ut spero venisti, videbis, ut soles, omnia quae intelleges nostra interesse, imprimis de Tullia mea, cuius de condicione quid mihi placeret scripsi ad Terentiam cum tu in Graecia esses; deinde de honore nostro. quod enim tu afuisti, vereor ut satis diligenter actum in senatu sit de litteris meis. illud praeterea μυστικώτερον ad te scribam, tu sagacius a odorabere. τῆσ δάμαρτόσ μου ὁ ἀπελεύθεροσ ( οἶσθα ὃν ) ἔδοξέ μοι πρώην, ἐξ ὧν ἀλογευόμενος παρεφθέγγετο, πεφυρακέναι τὰσ ψήφουσ ἐκ τῆσ ὠνῆσ τῶν ὑπαρχόντων τῶν τοῦ † νοήσῃς. † εἷσ δήπου . non queo tantum quantum vereor scribere; tu autem fac ut mihi tuae litterae volent obviae. haec festinans scripsi in itinere atque agmine. Piliae et puellae Caeciliae bellissimae salutem dices.
